\documentclass[../main.tex]{subfiles}
%\externaldocument[I-]{../files/Acronyms} 
% Conditionally load glossary setup for standalone compilation
    
\begin{document}
% Title page content
\section{Background}


Synthetic populations are artificially generated datasets that seek to replicate the characteristics of real study groups, including sociodemographic traits and behaviors. Their creation is often driven by the lack of access to detailed personal data due to unavailability, privacy and cost of sampling \citep{Chapuis2022GenerationPractices}. By excluding identifiable attributes while preserving realistic distributions, synthetic populations enable detailed modeling and analysis without privacy concerns \citep{Wu2018GeneratingDatasets}. These datasets are widely applied across disciplines, from modeling travel patterns and simulating disease spread to analyzing the effects of policy changes on diverse demographic groups \citep{Edwards2009TheSimObesity,Sallard2023TravelSwitzerland}. Additionally, they are integral to digital twin, which are dynamic representations of systems that combine synthetic populations with physical or virtual processes for more robust analysis and decision-making \citep{ElSaddik2018DigitalTechnologies}.

Despite their utility, generating realistic synthetic populations is challenging due to limited data availability and privacy requirements. Aggregated snapshots of population data often lack detail, while microdata, when available, is anonymized by reducing spatial granularity and generalizing attributes to broader categories \citep{Tanton2009OldAustralians}. For instance, in Australia's Census microdata, individual data is released at coarse spatial levels. For example, statistical area level 4 (SA4) regions of 100,000-300,000 people and attributes such as occupations are categorized into broad sectors (e.g., "health") \citep{AustralianBureauofStatistics2021CensusTools}. Although these aggregations protect privacy, they also create challenges for synthetic population generation, which must replicate observed totals and capture detailed population characteristics despite the loss of granularity due to available resources \citep{Sun2015ASynthesis}.

Synthetic population generation typically involves defining sociodemographic characteristics for synthetic individuals and households, assigning them to geographic zones, and, if needed, incorporating additional features such as social networks or activity schedules \citep{Somanath2025Activity-basedEquity}. Although existing methods focus on generating individual and household-level data \citep{Lenormand2013GeneratingMethods, Gargiulo2010AnHouseholds}, they often assume that spatial attributes are implicitly embedded in the data. Advanced approaches, such as enhancing microdata with neural networks, have improved representativeness \citep{Sun2015ASynthesis, Sun2018ASynthesis}, but they frequently fail to integrate diverse data sources or maintain household-level relationships. This leads to inconsistencies, such as generating individuals with household features but without corresponding members, which can bias equity assessments and household-level decision-making \citep{Qian2024ASynthesis}.

Efforts to address these limitations include matching algorithms to link individuals with home and workplace locations \citep{Djellab2020ImprovingInformation} and incorporating survey data, which provides behavioral insights but is often underutilized because of the assoicated financial costs and time conducted to conduct the survey \citep{Namazi-Rad2017AnSurveydata, Hawkins2023ASynthesis}. Deterministic reconstruction techniques, while useful, are constrained by outdated algorithms prone to issues such as cloning, sampling bias, and the zero-cell problem, where certain population groups expected to be included in the population are absent from the samples \citep{Acharya2022GenSyn:Sources}.

A flexible, scalable framework is therefore needed that integrates diverse data sources, uncovers latent variables in population samples, and bridges the gap between individual and household-level characteristics while maintaining logical consistency. Such a framework would enable more robust and adaptable synthetic population models, providing a foundation for addressing complex problems in fields ranging from urban planning to public health.

\section{Literature Review}
Methods for generating a synthetic population known as \gls{sr} can be categorized into three broad groups: 1. deterministic approaches, 2. \gls{co}, and 3. based on statistical learning and probabilistic simulation.

The characteristics of individuals in synthetic populations are primarily categorical. Continuous variables, such as income, are often grouped into defined ranges to facilitate their conversion into categorical data \citep{Borysov2019HowSynthesis}. This approach is necessary because generating population characteristics relies heavily on proportional fitting, as the large size of the data makes it challenging to directly capture overall population averages \citep{Wu2018GeneratingDatasets}. Typically, these characteristics are derived from microdata and adjusted to ensure consistency with aggregated control macrodata, such as averages and marginal distributions, for a given area.

One common approach involves replicating the same microdata individuals multiple times to match control totals. However, this method fails to produce people with diverse combination of characteristics present in the population \citep{Garrido2019PredictionModelling}. This approach can lead to entire groups of people being homogeneous. Lack of representation reduces heterogeneity present in the population. This issue, commonly referred to as the zero-cell problem, arises when samples from a known category are completely absent from the dataset, despite being known to exist in the actual population\citep{Sane2024ASynthesis}. This creates a paradigm in which the potential exists for generating individuals with specific combinations of attributes. However, the available pool of records is unable to fully capture associated characteristics, resulting in an incomplete representation of the population. To address this issue, Bayesian networks and probabilistic sampling techniques have been proposed to derive the missing categories (see \Cref{sec:bay} for a formal review).

\subsection{Deterministic Approaches}
\gls{sr} is a sample-based approach to create a synthetic population, with \gls{ipf} being the most common method due to its ability to match marginal distributions to known totals based on the sample source, provided individual microscopic data is available to represent individuals with specific characteristics. The \gls{sr} approach generates populations by randomly assigning individual characteristics, drawing values from observed distributions \citep{Barthelemy2012SyntheticSample, Gargiulo2010AnHouseholds}. However, \gls{ipf} has notable limitations, such as the zero-cell problem and its inability to simultaneously map individual and household population totals within a region \citep{Lovelace2014ALevels}.

To address the zero-cell problem, contingency tables can be smoothed by replacing "zeros" with small positive values (e.g., "pseudo-counts") before applying synthesis methods. This ensures that all cells are non-zero, allowing iterative methods such as \gls{ipf} to converge when structural zeros in the population are present \citep{Smith2017PopulationSampling,Lovelace2013TruncateMicrosimulation}. Another limitation of \gls{ipf} is that it is optimized based on fractional weights, creating a scaling issue because fractional individuals are not possible \citep{Namazi-Rad2017AnSurveydata}. This issue can be resolved through integerization techniques, such as quasi-integer rounding, to approximate integer populations \citep{Smith2017PopulationSampling}.

\gls{ipf}-based methods, while effective in replicating empirical data, are highly sensitive to data quality and sample size \citep{Gargiulo2010AnHouseholds}. These methods also enforce scaling to fit marginal constraints, which can lead to a cloning problem where identical records are repeatedly used to represent members of the population \citep{Whitworth2022SynthACS:Data}. Despite these limitations, \citet{Albiston2024ASynthesis} found that \gls{ipf} performs well when datasets include sufficiently detailed information about the population, such as well-populated contingency tables with minimal zero cells and diverse combinations of population characteristics. This is particularly true when targeting a limited number of characteristics for individuals in relatively large zones. However, \gls{ipf} cannot account for household relationships, making it inadequate for applications requiring joint imputation of individuals and households.

To address this limitation, \citet{Ye2009MethodologyPopulations} introduced \gls{ipu}, a method that simultaneously matches marginal totals for both households and individuals. \gls{ipu} produces a single joint weight that accounts for both persons and households, organizing the joint person-household sample data into a combined list. Unique cells in the resulting matrix represent household and person types. However, as the sample size and the number of characteristics increase, the matrix can become large and sparse, posing computational challenges and potentially preventing convergence \citep{Templ2017SimulationSimPop}.

\subsection{Combinatorial Optimization}
Another \gls{sr} method is \gls{co}, which constructs synthetic populations through a routine that selects records from available samples. A specified and customizable objective function evaluates the quality of the population, guiding the selection of sample combinations from the available data to define the resulting dataset.

The most commonly used \gls{co} algorithm is \gls{sa}, which constructs the population by minimizing the objective function through a metaheuristic search strategy. \gls{sa} iteratively draws from the sample and evaluates potential solutions, allowing stochastic adjustments that reduce differences in marginalized distributions for each modeled characteristic, guiding the process toward an optimized solution \citep{Abraham2012PopulationLevels, Whitworth2022SynthACS:Data, Kim2016AModel, Duran-Heras2018ComparisonMethod}.

One advantage of \gls{co} over \gls{ipf} and \gls{ipu} is its ability to avoid rounding issues when matching marginal distributions. This ensures a feasible solution even when distributions are not perfectly aligned \citep{Duran-Heras2018ComparisonMethod}. Through its optimization process, \gls{co} can refine marginal distributions and address various aspects of solution quality using a customized fitness function \citep{Huang2001AMicrodata.}. Additionally, while cloning is required when the sample is smaller than the population, \gls{co} offers the flexibility to constrain the search to meet heterogeneity requirements, such as preventing the over-cloning of specific samples \citep{Zhuge2017ATransportation, Ma2015SyntheticValidations, Harland2012CreatingTechniques}. Furthermore, \gls{co} overcomes issues such as zero-cell and sample bias problems, which are common in \gls{ipf} \citep{Kim2016AModel}. However, because the algorithm is dependent on the sample, challenges such as cloning and the relative simplicity of alternative methods have limited the widespread adoption of \gls{co} \citep{Wang2021ImprovedPopulations}.

% While most algorithms primarily focus on matching marginal distributions, \citet{Zhuge2017ATransportation} introduced a heuristic method, Pop-H, which leverages validation objectives to improve the standard deviation of generated characteristics, resulting in more representative synthetic populations. Building on this, \citet{Bigi2024SyntheticMobility} emphasized the importance of distributional assumptions during the generation process, highlighting their role in enhancing the accuracy and validity of synthetic population modeling.

\subsection{Statistical Learning} \label{sec:bay}
Probabilistic learning approaches have been proposed as alternatives to the deterministic methods used for synthetic reconstruction \citep{Farooq2013SimulationSynthesis, Ilahi2019IntegratingJakarta, Sun2015ASynthesis, Zhang2019ConnectedSimulation}. This is because microsamples are not always available and statistical approaches are more scalable as the number of characteristics increases \citep{Borysov2019HowSynthesis}. 

While some studies opted to use Markov Chains for probabilistic population synthesis \citep{Farooq2013SimulationSynthesis}, others adopted data-driven inferential methods such as Bayesian Networks \citep{Sun2015ASynthesis,Zhang2019ConnectedSimulation} or graph based representations through \glspl{gan} \citep{Lederrey2022DATGAN:Data}. 

Markov Chains, the simplest of these methods, can capture correlations among variables sequentially, which can be challenging to model when the sequence, or ordering, of variables is unknown and complex interdependencies exist among population characteristics \citep{Saadi2016HiddenSynthesis}. The most basic sampling technique is the Monte Carlo method  that uses raw input data regarding the distribution of characteristics. This method is labeled a direct sampling, in which the process relies on raw data. Individual characteristics can either be drawn sequentially or several at a time if the input data includes records on relationships among them. However, with an increasing number of population characteristics, multi-way joint distributions of characteristics together become infeasible. Therefore, a realistic dependency structure between individual characteristics to household must be determined in advance \citep{Sun2018ASynthesis}. 

Subsequently, a synthetic population can be generated iteratively using Bayesian rules on conditional distributions \citep{Gargiulo2010AnHouseholds}. \Gls{hs} represents the most fundamental model for this type of sampling method \citep{Barthelemy2012SyntheticSample}. An initial decision is needed to establish the hierarchy. This requires specifying which characteristic to sample first, then determining subsequent characteristics conditioned on the previously selected values. This process continues iteratively until the final characteristic is sampled.
An alternative to hierarchical sampling is the use of a \gls{bn} for synthetic population generation, which extends and generalizes the principles of \gls{hs} through a graphical model \citep{Sun2015ASynthesis,Ye2017PopulationSamples}. Network-based samplers operate iteratively, beginning with a root node (or nodes) that represent initial characteristics and then traversing the network's edges to systematically sample additional connected characteristics according to predefined sampling criteria. These samplers offer a significant advantage over \gls{hs} by leveraging classical learning techniques to automatically determine the graphical model and/or estimate missing parameters, while adhering to rules and constraints derived from scattered statistics \citep{Thiriot2008GenerateStatistics}.

Furthermore, if known distributions of characteristics are available, a graphical model can be directly derived from them, meaning that sample data is not required to achieve results. However, when initialized with microdata, the model should be estimated using conditionals derived from the microdata \citep{Sun2015ASynthesis}. 

Although Bayesian network models provide access to aggregated characteristics across different geographical zones, they can only precisely match the distribution of a single characteristic (the mother node). According to the conditional distribution theorem, once the distribution of the mother node and a conditional probability table are established, the distribution of child nodes is uniquely defined \citep{Jafari2022BuildingModels}. This implies that it is impossible to simultaneously satisfy both the marginal and conditional distributions for all characteristics. A practical solution is to keep the Bayesian network as simple as possible - capturing only statistically significant dependencies and allowing the unknown features to vary \citep{Zhang2019ConnectedSimulation,Aemmer2022GenerativeCharacteristics}.

Generative models extend the same programme to high-dimensional tables. \citet{Lederrey2022DATGAN:Data} condition a \gls{gan} on a directed acyclic graph of attribute dependencies, so that sampled records follow a learned association structure rather than being drawn one attribute at a time. The limitation that matters for this paper is not statistical but structural. A generative model is fitted to the \emph{observed joint distribution}, whereas the rules that make a record logically possible are known a priori, from the census dictionary, and they constrain combinations the microsample may never have contained often enough to move a loss function. Nothing in an adversarial or likelihood objective penalises a record that gives an occupation to a nine-year-old, or a dwelling that both rents and carries a mortgage: such records are rare in the training sample rather than impossible in the world. Generated records are therefore checked against an explicit rule set in published applications, and those that fail are repaired or rejected \citep{Kim2023ASynthesis}, so the constraint layer becomes a separate stage whose function is to undo what the generator produced.

It is argued here that this ordering is the wrong way round. Once the structural rules have to be enumerated rule by rule and applied record by record---and they do, in every application that reports a repair rate---the structural knowledge is already doing the substantive work, and the generator is contributing only the association structure, which the observed conditional distributions also contain. If the enumeration of allowable combinations is required in any case, it is more economical and more defensible to construct candidate records \emph{inside} the constraint set from the outset, so that feasibility is the method rather than a clean-up applied to its output. That is the position adopted here: candidate records are built by constraint-guided resampling from the census conditionals, and the pool is expanded by a mixed-integer program that admits only rule-feasible households. The approach is not offered as a refutation of generative synthesis in general, but as the appropriate choice when the control table, not the microsample, is the binding object---which is the case in the zone-level fitting problem addressed below.


%While deep learning methods provide promising approaches for population synthesis, long-standing issues of machine learning such as the lack of interpretability (i.e., the black-box nature of machine learning models) and the tendency to overfit the training data remain viable concerns, however, these approaches are able to identify the relationships and the generated result does in deep retain an individuals attributes.

% An alternative to sample-based approaches is the use of sample-free methods, which have been discussed extensively in the population synthesis literature for cases where a representative sample is unavailable \citep{Gargiulo2010AnHouseholds,Lenormand2013GeneratingMethods}. Unlike sample-based methods, sample-free approaches do not rely on a sample record file to construct a synthetic population. Instead, they used heuristics to better preserve the geographical heterogeneity of the resulting synthetic population. However, one major limitation of sample-free approaches is the lack of a clearly identified source of variability in the simulated synthetic population. Additionally, these methods must address inconsistencies between margins extracted from different cross-tabulations. Setting multidimensional weights for population characteristics without reliable seed data to act as a source of variability further complicates the process and requires extensive computational resources \citep{Namazi-Rad2014GeneratingDynamics}. In contrast, sample-based methods calculate these weights directly from the sample data.

\subsection{Hybrid Methods}

Hybrid approaches have been developed to address the limitations of sample-based and statistical learning methods by combining the strengths of both. These approaches expand or construct a sample pool using statistical learning methods, which employ fitting algorithms like \gls{ipu} or \gls{sa} \citep{Sun2015ASynthesis, Tozluoglu2023APatterns}. For example, \citet{Fournier2021IntegratedMethod} and \citet{Casati2015SyntheticRaking} utilized a \gls{mcmc} algorithm within a proportional fitting framework. This hybrid method enhances the quality of the microdata pool by increasing sample heterogeneity and subsequently applying a proportional fitting method. As a result, it achieves better alignment with marginal controls compared to traditional methods like \gls{ipf} alone, as it uncovers more latent variables through the statistical learning process \citep{Muller2017ASYNTHESIS}. Furthermore, it reduces the reliance on expensive and detailed data collection efforts by providing an alternative fitting approach capable of overcoming the zero-cell problem \citep{Zhou2022CreatingMicrosimulation}.

 In addition, advances in generating realistic populations for regions with limited detailed information have been further expanded through the use of alternative deep learning models and various proposed methods \citep{Jutras-Dube2024Copula-basedGeneration}. For instance, \citet{Borysov2019HowSynthesis} and \citet{Aemmer2022GenerativeCharacteristics} utilized Variational Autoencoders to expand micro-agent pools, creating new samples with latent variables derived from the sampled population. Similarly, \citet{Ilahi2019IntegratingJakarta} incorporated Bayesian networks, while \citet{Saadi2018AnFusion} employed Hidden Markov Models (HMM). Bayesian network-based synthesis has been further refined by methods such as generalized raking, as demonstrated by \citet{Ilahi2019IntegratingJakarta} and \citet{Zhou2022CreatingMicrosimulation}, who reported a good fit to marginal controls.

Additionally, hierarchical mixture models have been used to characterize the joint distribution of individual and household attributes. This framework has shown significant potential for generating representative household structures, effectively reproducing the statistical distributions of real populations \citep{Sun2018ASynthesis}. Despite these advancements, no approaches have yet leveraged \gls{co} as a sampler, with most relying on \gls{ipu} or \gls{ipf}.

\section{Proposed Framework}

What is described here is an \emph{integration and reproducibility} study of a construction-first pipeline. Each element described below---constraint-guided resampling from observed conditionals, mixed-integer selection, household reconstruction, and \gls{ipu}/\gls{sa} fitting---is individually established in the literature; the object of study is their combination, arranged so that feasibility is established at construction rather than repaired afterwards, and measured through a controlled comparison of what each component contributes. The framework is designed to account for individual heterogeneity, preserve household relationships, and align a synthetic population with a large number of zone-level marginal control tables.

The framework consists of five stages:
\begin{itemize}
\item \textbf{Stage 1 --- Feasible pool construction}: Candidate records are drawn from the census microdata and from structurally feasible but unobserved cells reconstructed by constraint-guided conditional resampling. The 11 census-consistency rules are expressed as explicit predicates over the attribute domains, so the pool is feasible by construction, cell-complete over the modelled attributes, deterministic given a seed, and requires no post-hoc repair pass.

\item \textbf{Stage 2 --- Pool expansion}: Where a zone's control tables require a cell that the pool cannot supply at all, a mixed-integer program selects the smallest set of rule-feasible donor--shell households that restores coverage, subject to every supplied marginal table simultaneously. Coverage is soft, so a cell that no donor can source is reported rather than invented.

\item \textbf{Stage 3 --- Household Relationship Reconstruction}: The pool is processed with a greedy matching algorithm to reconstruct household relationships while maintaining consistency across individuals, and the records added by expansion enter as whole households rather than as free-floating persons.

\item \textbf{Stage 4 --- Adjustment to Marginal Distributions}: Households drawn from the pool are selected to match zone-level marginal control distributions using \gls{ipu}, \gls{sa}, or a tolerance-relaxed mixed-integer fit. This stage captures both spatial and attribute-level heterogeneity.

\item \textbf{Stage 5 --- Validation}: The synthetic population is validated against held-out census tables that were not used in fitting.
\end{itemize}

Through these steps a synthetic population is generated that supports simultaneous imputation of individual and household attributes, integerization of the population, synthesis of zero-cells, and integration of diverse control variables. These properties are not claimed to be achieved ``better'' than in all prior work; what is demonstrated, with released code and a single reproducible run, is that they are achieved jointly and measurably.



Cloning---the repeated reuse of a small set of records across zones---has historically limited the adoption of \gls{sa} for population synthesis \citep{Huang2001AMicrodata.}. Widening the candidate pool reduces, though does not eliminate, this cloning: expansion adds supported-but-unobserved records, and simulated annealing then selects a subset of households from the widened pool so that the synthetic population matches the aggregate totals of the real population. The reduction is a property of the \emph{larger pool}, and residual cloning is structural when a finite pool is reused across many zones.

\gls{ipu} is included as a conventional baseline for comparison. Although \gls{ipu} is widely used, it is known to suffer from cloning and difficulties in meeting household-level constraints.

The design is therefore a two-factor comparison. The \emph{pool} is either the raw microdata or an explicitly feasible constructed pool, with and without mixed-integer expansion; the \emph{fitter} is \gls{sa}, \gls{ipu}, or a mixed-integer fit. A detailed evaluation of the resulting seven strategies---SA-Only, IPU-Only, CP-pool+SA, CP-pool+IPU, CP-MIP, and \gls{sa} and \gls{ipu} applied to the MIP-expanded pool---focusing on the validity of the produced populations, adherence to constraints, and the measured effect of each pipeline component, is provided in the sections that follow.

\ifstandalone
    % This block is executed only when the subfile is compiled standalone
    \bibliographystyle{elsarticle-harv}
    \bibliography{references}
\fi
\end{document}